\documentclass[11pt, headheight=30pt]{scrartcl}
\usepackage[white, nologo, hw]{darkWhite}
\usepackage{mathrsfs}

\newcommand{\BBB}{\mathscr{B}}
\DeclareMathOperator{\cov}{cov}
\DeclareMathOperator{\var}{var}

\title{18.675 Problem Set 6}
\begin{document}
% PROBLEMS SET 6 FOR 18.675 (FALL 2023)
% Problem 1. Let X be a martingale (resp. supermartingale) on some filtered probability space, and
% let T be an a.s. finte stopping time. Prove that \EE[X_T ] = \EE[X_0] (resp. \EE[X_T ] ≤ \EE[X_0]) if either one
% of the following conditions holds:
% (i) The process X is bounded (∃M > 0 : ∀n ≥ 0, |Xn| ≤ M a.s.).
% (ii) The process X has bounded increments (∃M > 0 : ∀n ≥ 0, |Xn+1 − Xn| ≤ M a.s.) and
% \EE[T] < ∞.
\section{}
\begin{prob*}
    Let $X$ be a martingale (resp. supermartingale) on some filtered probability space, and
    let $T$ be an a.s. finte stopping time. Prove that $\EE[X_T ] = \EE[X_0]$ (resp. $\EE[X_T ] \leq \EE[X_0]$) if either one
    of the following conditions holds:
    \begin{enumerate}
        \item The process $X$ is bounded ($\exists M > 0 : \forall n \geq 0, |X_n| \leq M$ a.s.).
        \item The process $X$ has bounded increments ($\exists M > 0 : \forall n \geq 0, |X_{n+1} - X_n| \leq M$ a.s.) and
        $\EE[T] < \infty$.
    \end{enumerate}
\end{prob*}

The supermartingale case immediately implies the martingale case,
so we only prove the supermartingale case.
First off, if $T$ were bounded, then the Bounded Optional Stopping Theorem
would conclude. In particular, $\EE[X_{T\land n}] \leq \EE[X_0]$ for any
constant $n\geq 0$. So all we want to show is that $\EE[X_{T\land n}] \to \EE[X_T]$.
All we have to do is rip through \vocab{Fatou-Lebesgue} by providing an integrable
function which dominates all $X_{T\land n}$ almost surely.

In the first case, $X_{T\land n}$ is dominated by $M$ almost surely, so we conclude.

In the second case, we need to use the condition $\EE[T] < \infty$ by expanding
in terms of increments. We know
\[
    X_{T\land n} = X_0 + \sum_{t = 0}^{n-1} (X_{t+1} - X_t)\id_{t\leq T}
\]
thus
\begin{align*}
    \abs{X_{T\land n}} &\leq \abs{X_0} + \sum_{t\geq 0} \abs{X_{t+1} - X_t}\id_{t\leq T}\\
    &\leq \abs{X_0} + M\sum_{t\geq 0} \id_{t\leq T}\\
    &\leq \abs{X_0} + M\EE[T] < \infty
\end{align*}
The RHS of the top line is independent of $n$, and is our dominating function.
We conclude.


% Problem 2. Let T be an (Fn, n ≥ 0)-stopping time such that for some integer N > 0 and  > 0,
% \PP[T ≤ N + n | Fn] ≥ , for every n ≥ 0.
% Show that \EE[T] < ∞.
\section{}
\begin{prob*}
    Let $T$ be an $(\mathscr{F}_n, n \geq 0)$-stopping time such that for some integer $N > 0$ and $\epsilon > 0$,
    $\PP[T \leq N + n | \mathscr{F}_n] \geq \epsilon$, for every $n \geq 0$.
    Show that $\EE[T] < \infty$.
\end{prob*}
By definition, the condition says that for all measurable $A\in \mathscr{F}_n$,
\[
    \PP[A \cap \{T \leq N + n\}] \geq \epsilon \PP[A].
\]
Sure, let's pick $A = \{T > n\}$! Then,
\[
    \PP[n < T \leq N + n] \geq \epsilon \PP[T > n].
\]
Then,
\begin{align*}
    \EE[T] &= \sum_{n\geq 0} \PP[T > n]\\
    &\leq \frac{1}{\eps}\sum_{n\geq 0} \PP[n < T \leq N + n]\\
    &\leq \frac{1}{\eps}\sum_{0\leq i < N} \sum_{n\geq 0} \PP[n\cdot N + i < T \leq (n+1)\cdot N + i]\\
    &\leq \frac{1}{\eps}\sum_{0\leq i < N} \PP[i < T]\\
    &\leq \frac{1}{\eps}\sum_{0\leq i < N} 1 < \infty
\end{align*}

% Problem 3. Let (Xn, n ≥ 0) be a sequence of [0, 1]-valued random variables, which satisfy
% the following property. First, X_0 = a a.s. for some a ∈ (0, 1) and for n ≥ 0,
% P
% 
% Xn+1 =
% Xn
% 2
% | Fn
% 
% = 1 − Xn = 1 − P
% 
% Xn+1 =
% Xn + 1
% 2
% | Fn
% 
% ,
% where Fn = σ(Xk, 0 ≤ k ≤ n). Here, we have denoted \PP[A | G] = \EE[1A | G].
% (i) Prove that (Xn, n ≥ 0) is a martingale that converges in L
% p
% for every p > 1.
% (ii) Check that \EE[(Xn+1 − Xn)
% 2
% ] = \EE[Xn(1 − Xn]/4. Then determine \EE[X∞(1 − X∞)] and deduce
% the law of X∞.
\section{}
\begin{prob*}
    Let $(X_n, n \geq 0)$ be a sequence of $[0, 1]$-valued random variables, which satisfy
    the following property. First, $X_0 = a$ a.s. for some $a \in (0, 1)$ and for $n \geq 0$,
    \[
        P
        \left\{
            X_{n+1} = \frac{X_n}{2} \bigg| \mathscr{F}_n
        \right\}
        = 1 - X_n = 1 - P
        \left\{
            X_{n+1} = \frac{X_n + 1}{2} \bigg| \mathscr{F}_n
        \right\},
    \]
    where $\mathscr{F}_n = \sigma(X_k, 0 \leq k \leq n)$. Here, we have denoted $\PP[A | G] = \EE[1_A | G]$.
    \begin{enumerate}
        \item Prove that $(X_n, n \geq 0)$ is a martingale that converges in $L^p$ for every $p > 1$.
        \item Check that $\EE[(X_{n+1} - X_n)^2] = \EE[X_n(1 - X_n)]/4$. Then determine $\EE[X_\infty(1 - X_\infty)]$ and deduce
        the law of $X_\infty$.
    \end{enumerate}
\end{prob*}

For the first part, by Doob's Martingale Convergence Theorem, it suffices to show that
$X_n$ is bounded in $L^p$ for every $p > 1$. But this is clear, since $X_n \in [0, 1]$,
so $\EE[\abs{X_n}^p] \leq 1$ uniformly in $n$. As a bonus, we know that this limit $X_\infty$
is also an almost-sure limit, and $X_n = \EE[X_\infty | \mathscr{F}_n]$ almost surely for all $n\geq 0$.

For the second part, we compute
\begin{align*}
    \EE[(X_{n+1} - X_n)^2] &= \EE[ \EE[(X_{n+1} - X_n)^2 | \mathscr{F}_n]]\\
    &= \EE\left[X_n\left(\frac{X_n + 1}{2} - X_n\right)^2 + (1 - X_n)\left(\frac{X_n}{2} - X_n\right)^2\right]\\
    &= \frac14 \EE[X_n(1 - X_n)^2 + (1-X_n)X_n^2]\\
    &= \frac14 \EE[X_n(1 - X_n)].
\end{align*}

Now, we know that $X_n \to X_\infty$ almost surely,
so by the Dominated Convergence Theorem,
\[
    \EE[X_\infty(1 - X_\infty)] = \lim_{n\to \infty} \EE[X_n(1 - X_n)] = 4 \lim_{n\to \infty} \EE[(X_{n+1} - X_n)^2] = 0
\]
where the last equality is because $X_n$ converge in $L^2$.

However, $\EE[X_\infty(1-X_\infty)]$ is nonnegative on $[0,1]$, with
equality at $\{0,1\}$ only. Thus, $X_\infty$ is almost surely $0$ or $1$.
$X$ is also UI, thus by optional stopping theorem, $\EE[X_\infty] = \EE[X_0] = a$.
Thus, $\PP(X_\infty = 1) = a$ and $\PP(X_\infty = 0) = 1 - a$.


% Problem 4. Let (Xn, n ≥ 0) be a sequence of independent and identically distributed real-valued
% integrable random variables. We let Sn = X1 + · · · + Xn (with S0 = 0) be the associated random
% walk and T an (Fn)-stopping time, where Fn = σ(Xk, k ≤ n).
% (1) Show that if the random variables Xi are non-negative, then
% \EE[ST ] = \EE[T]\EE[X1].
% (2) Show that if \EE[T] < ∞, then
% \EE[ST ] = \EE[T]\EE[X1].
% (3) Suppose that \EE[X1] = 0 and set Ta = inf{n ≥ 0 : Sn ≥ a}, for some a > 0. Show that
% \EE[Ta] = ∞.
% (4) Suppose that \PP[X1 = +1] = 2
% 3 = 1 − \PP[X1 = −1] and set Ta = inf{n ≥ 0 : Sn ≥ a}, for some
% a > 0. Find \EE[Ta]. (You cannot assume that \EE[Ta] < ∞.).
\section{}
Let $(X_n, n \geq 0)$ be a sequence of independent and identically distributed real-valued
integrable random variables. We let $S_n = X_1 + \cdots + X_n$ (with $S_0 = 0$) be the associated random
walk and $T$ an $(\mathscr{F}_n)$-stopping time, where $\mathscr{F}_n = \sigma(X_k, k \leq n)$.
% \begin{enumerate}
%     \item Show that if the random variables $X_i$ are non-negative, then
%     $\EE[S_T] = \EE[T]\EE[X_1]$.
%     \item Show that if $\EE[T] < \infty$, then
%     $\EE[S_T] = \EE[T]\EE[X_1]$.
%     \item Suppose that $\EE[X_1] = 0$ and set $T_a = \inf\{n \geq 0 : S_n \geq a\}$, for some $a > 0$. Show that
%     $\EE[T_a] = \infty$.
%     \item Suppose that $\PP[X_1 = +1] = 2/3 = 1 - \PP[X_1 = -1]$ and set $T_a = \inf\{n \geq 0 : S_n \geq a\}$, for some
%     $a > 0$. Find $\EE[T_a]$. (You cannot assume that $\EE[T_a] < \infty$.).
% \end{enumerate}
\begin{prob}
    Show that if the random variables $X_i$ are non-negative, then
    $\EE[S_T] = \EE[T]\EE[X_1]$.
\end{prob}

\idea{
    [Main idea]
    This problem is very interesting! The key idea is that $\EE[X_i\id_{T = n}]$ does
    \emph{not} have to equal $\EE[X_i]\PP(T = n)$, because $T$ is not necessarily
    independent of $X_i$. However, we notice that postponing $T$ does not change the
    fact that we wil \emph{eventually} count $X_i$; symbolically, we really care
    about $\EE[X_i\id_{T\geq i}]$. Flipping this to $\EE[X_i\id_{T < i}]$, we
    observe that $\{T < i\}$ and $X_i$ really are independent, so we can compute
    everything.
}


If $\EE[X_1] = 0$, then $X_1 = 0$ almost surely and the problem is a bit stupid;
assume otherwise. Now, by monotone convergence or whatever,
\[
    \EE[S_T] = \sum_{n = 0}^\infty \EE[S_n \id_{T = n}] + \EE[S_\infty \id_{T = \infty}].
\]
If $\PP(T = \infty) > 0$, then $\EE[X_1]\EE[T] = \infty$, and the LHS is also $\infty$
(because $S_\infty$ is also infinite almost surely, if you interpret this to mean
the almost-sure limit; this guy might also be undefined so whatever). So we can assume
that $\PP(T = \infty) = 0$. Then, we observe the LHS is just
\[
    \sum_{n = 0}^\infty \sum_{i\leq n} \EE[X_i\id_{T = n}]
\]
By Fubini's theorem we can interchange the sums.
\[
    \sum_{i = 0}^\infty \sum_{n\geq i} \EE[X_i\id_{T = n}]
\]
Now, we observe that $\PP(T \geq i) = 1 - \PP(T \leq i - 1)$ which lives in $\mathscr{F}_{i-1}$
and is thus independent of $X_i$. Thus, we can write
\[
    = \sum_{i = 0}^\infty \EE[X_i]\PP(T \geq i) = \EE[X_1]\EE[T]
\]
as desired.

\begin{prob}
    Show that if $\EE[T] < \infty$, then
    $\EE[S_T] = \EE[T]\EE[X_1]$.
\end{prob}

First off the fact that $\EE[T]$ exists means that $\PP(T = \infty) = 0$.

Now, we repeat the above proof with $[0, \infty]$ replaced by $\RR$.
Consider the partial sums
\[
    \sum_{n = 0}^m S_n\id_{T = n} \stackrel{m\to \infty}{\longrightarrow} S_T   
\]
which converge almost surely. We need to show the expectation of the absolute
value of the LHS is bounded so we can apply Dominated Convergence.

Of course, this is just the previous problem, because
\[
    \EE\left[\abs{\sum_{n = 0}^m S_n\id_{T = n}}\right]\leq \sum_{n = 0}^\infty \EE[\tilde{S}_n\id_{T = n}]\leq \EE[T] \EE[\abs{X_1}] < \infty
\]
where $\tilde{S} = \sum_{i = 0}^n \abs{X_i}$; note that $\abs{X_i}$ is integrable if $X_i$
is.

Great! So we really care about the sum
\begin{align*}
    \sum_{n = 0}^\infty \EE[S_n\id_{T = n}]
    &= \sum_{n = 0}^\infty \sum_{i\leq n} \EE[X_i\id_{T = n}]\\
    &= \sum_{i = 0}^\infty \sum_{n\geq i} \EE[X_i\id_{T = n}]\\
    &= \sum_{i = 0}^\infty \EE[X_i]\PP(T \geq i)\\
    &= \EE[X_1]\EE[T]
\end{align*}
by the same arguments as the first time; use Tonelli rather than Fubini.
\begin{prob}
    Suppose that $\EE[X_1] = 0$ and set $T_a = \inf\{n \geq 0 : S_n \geq a\}$, for some $a > 0$. Show that
    $\EE[T_a] = \infty$.
\end{prob}

Applying the previous problem to $T = T_a$, we observe that if $\EE[T_a] < \infty$
(such that $T_a < \infty$ almost surely) then $\EE[S_T] = \EE[T_a]\EE[X_1] = 0$.
However, in this case $S_T \geq a$ almost surely thus $\EE[S_T]\geq a$, contradiction.

Thus $\EE[T_a] = \infty$.

\begin{prob}
    Suppose that $\PP[X_1 = +1] = 2/3 = 1 - \PP[X_1 = -1]$ and set $T_a = \inf\{n \geq 0 : S_n \geq a\}$, for some
    $a > 0$. Find $\EE[T_a]$. (You cannot assume that $\EE[T_a] < \infty$.).
\end{prob}

Well, if we knew that $\EE[T_a] < \infty$,
then $\EE[S_T] = \EE[T]\EE[X_1] = \frac13 \EE[T]$; of course,
in this case, $\EE[S_T] = \ceil{a}$,
so we would find $\EE[T] = 3\ceil{a}$.

Okay, let's just compute an upper bound on $\EE[T_a]$!
The probability that $T_a$ occurs after timestep $n\gg a$ is the probability
that at least $\frac12 n - \frac12 a$ of the $X_i$ are $-1$. This binomial sum is
\[
    \sum_{k = \frac12 n -\frac12 a}^n \binom{n}{k}
    \left(\frac13\right)^k \left(\frac23\right)^{n-k} 
\]
Now, the claim is that individual terms $\binom{n}{k}
\left(\frac13\right)^k \left(\frac23\right)^{n-k}$ are decreasing; this is true
for sufficiently large $n$, because the ratio of successive $\binom{n}{k}$ increases
by at most $1 + O(\frac{a}{n})$ while the ratio of successive $\left(\frac13\right)^k \left(\frac23\right)^{n-k}$
is exactly $1/2$.

Thus, this entire sum is bounded above by
\begin{align*}
    \left(\frac12 n + \frac12 a\right) \binom{n}{\frac12 n - \frac12 a} \left(\frac13\right)^{\frac12 n - \frac12 a} \left(\frac23\right)^{\frac12 n + \frac12 a}
    \leq 2^n \cdot \left(\frac29\right)^{n/2}  = \left(\frac{2\sqrt{2}}{3}\right)^n 
\end{align*}
where I have dropped subleading terms in $a$.

Okay, fine! Now, recalling $\EE[T] = \sum_{n = 0}^\infty \PP(T > n)$, we can truncate
a finite prefix and bound the remaining terms by
\[
    \sum_{n = k}^\infty \left(\frac{2\sqrt{2}}{3}\right)^n < \infty
\]
This concludes.

Alternatively, say something about a Chernoff bound to kill the problem.


% Problem 5. Suppose that X1, X2, · · · are independent random variables with
% \PP[X = +1] = p, \PP[X = −1] = q,
% where p ∈ (0, 1), q = 1 − p and p 6= q. Suppose that a and b are integers with 0 < a < b. Define
% Sn = a + X1 + · · · + Xn, T = inf{n : Sn = a or Sn = b}.
% 1
% 2 PROBLEMS SET 6 FOR 18.675 (FALL 2023)
% Let Fn = σ(X1, · · · , Xn). Prove that
% Mn =
% 
% q
% p
% Sn
% and Nn = Sn − n(p − q)
% define martingales M and N. Deduce the values of \PP[ST = 0] and \EE[T].
\section{}
\setcounter{prob*}{4}
\begin{prob*}
    Suppose that $X_1, X_2, \cdots$ are independent random variables with
    $\PP[X = +1] = p, \PP[X = -1] = q$,
    where $p \in (0, 1)$, $q = 1 - p$ and $p \neq q$. Suppose that $a$ and $b$ are integers with $0 < a < b$. Define
    $S_n = a + X_1 + \cdots + X_n$, $T = \inf\{n : S_n = 0 \text{ or } S_n = b\}$.
    Let $\mathscr{F}_n = \sigma(X_1, \cdots , X_n)$. Prove that
    \[
        M_n =
        \left(
            \frac{q}{p}
        \right)^{S_n}
        \qquad
        \text{and}
        \qquad
        N_n = S_n - n(p - q)
    \]
    define martingales $M$ and $N$. Deduce the values of $\PP[S_T = 0]$ and $\EE[T]$.
\end{prob*}

First, we show $M$ and $N$ are martingales. Indeed,
\begin{align*}
    \EE[M_{n+1} | \mathscr{F}_n] &= \EE\left[
        \left(
            \frac{q}{p}
        \right)^{S_{n+1}}
        \bigg|
        \mathscr{F}_n
    \right]\\
    &= \EE\left[
        \left(
            \frac{q}{p}
        \right)^{S_n + X_{n+1}}
        \bigg|
        \mathscr{F}_n
    \right]\\
    &= \left(\frac{q}{p}\right)^{S_n}
    \EE\left[
        \left(
            \frac{q}{p}
        \right)^{X_{n+1}}
        \bigg|
        \mathscr{F}_n
    \right]\\
    &= \left(\frac{q}{p}\right)^{S_n}
    \left( \frac{p}{q}\cdot \PP[X_{n+1} = -1] + \frac{q}{p}\cdot \PP[X_{n+1} = +1]\right)\\
    &= M_n
    \left( \frac{p}{q}\cdot q + \frac{q}{p}\cdot p\right)\\
    &= M_n
\end{align*}
and
\begin{align*}
    \EE[N_{n+1} | \mathscr{F}_n] &= \EE\left[
        S_{n+1} - (n+1)(p - q)
        \bigg|
        \mathscr{F}_n
    \right]\\
    &= \EE\left[
        S_n + X_{n+1} - (n+1)(p - q)
        \bigg|
        \mathscr{F}_n
    \right]\\
    &= S_n - n(p - q) + \EE[X_{n+1}] - (p - q)\\
    &= N_n
\end{align*}


For $M$, observe that $S_{T\land n}$ is a bounded martingale thus optional
stopping theorem applies and $\EE[M_T] = \EE[M_0] = \left(\frac{q}{p}\right)^a$. For $N$, by arguments identical
to $4$d we know $\EE[T] < \infty$, plus the increments are bounded thus
optional stopping theorem applies and $\EE[N_T] = \EE[N_0] = a$.

On the other hand, we know that $S_T\in \{0, b\}$ almost surely. Thus we can extract
the probabilities:
\[
    \PP(S_T = 0)\cdot \left(\frac{q}{p}\right)^0 + \PP(S_T = b)\cdot \left(\frac{q}{p}\right)^b = \left(\frac{q}{p}\right)^a    
\]
and
\[
    \PP(S_T = 0) + \PP(S_T = b) = 1
\]
thus
\[
    \PP(S_T = 0) + (1 - \PP(S_T = 0))\cdot \left(\frac{q}{p}\right)^b = \left(\frac{q}{p}\right)^a    
\]
and
\[
    \PP(S_T = 0) = \frac{\left(\frac{q}{p}\right)^a - \left(\frac{q}{p}\right)^b}{1 - \left(\frac{q}{p}\right)^b}.
\]
We also know
\[
    \PP(S_T = 0)\cdot 0 + \PP(S_T = b)\cdot b - \EE[T](p-q) = a  
\]
thus
\[
    \EE[T] = \frac{\PP(S_T = b)\cdot b + a}{p - q}    
\]



% Problem 6. Let f : [0, 1] 7→ R be Lipschitz, that is, suppose that, for some K < ∞ and all
% x, y ∈ [0, 1]
% |f(x) − f(y)| ≤ K|x − y|.
% Denote by fn the simplest piecewise linear function agreeing with f on {k2
% −n
% : k = 0, 1, · · · , 2
% n}.
% Set Mn = f
% 0
% n
% . Show that Mn converges a.s. and in L
% 1 and deduce that f is the indefinite integral of
% a bounded function.
\section{}
\begin{prob*}
    Let $f : [0, 1] \to \RR$ be Lipschitz, that is, suppose that, for some $K < \infty$ and all
    $x, y \in [0, 1]$
    \[
        |f(x) - f(y)| \leq K|x - y|.
    \]
    Denote by $f_n$ the simplest piecewise linear function agreeing with $f$ on $\{k2^{-n} : k = 0, 1, \cdots , 2^n\}$.
    Set $M_n = f_n'$. Show that $M_n$ converges a.s. and in $L^1$ and deduce that $f$ is the indefinite integral of
    a bounded function.
\end{prob*}
To be explicit, on the interval $[k\cdot 2^{-n}, (k + 1)\cdot 2^{-n}]$, $f_n$
is linear and agrees with $f$ at the endpoints.

% By the Lipschitz condition,
% \[
%     \abs{f(x) - f(k\cdot 2^{-n})}\leq K\cdot 2^{-n}    
% \]
% and
% \[
%     \abs{f(x\cdot 2^{-n}) - M_n(x)}\leq K\cdot 2^{-n}
% \]
% thus
% \[
%     \abs{f(x) - f_n(x)}\leq 2K\cdot 2^{-n}    
% \]
% and $f_n$ converges uniformly to $f$. This immediately implies $M_n$

We claim that there is a Martingale hidden in here. Indeed, let $\Omega = [0,1]$;
let our filtration be
\[
    \mathfrak{F}_n = \sigma(\{[k2^{-n}, (k + 1) 2^{-n}] : k = 0, 1, \cdots , 2^{n-1}\})
\]
Then, $M_n$ is $\mathfrak{F}_n$-measurable, and
\[
    \EE[M_{n+1} | \mathfrak{F}_n] = \EE[f_{n+1}' | \mathfrak{F}_n] = f_n'    
\]
(this is true because it is true on each segment $[k2^{-n}, (k+1)2^{-n}]$).

Thus $M_n$ is a martingale! But it is also bounded by $\pm K$, because $f$ is Lipschitz.
Thus $M_n$ is $L^1$ bounded and uniformly integrable.
Thus, by Doob's Martingale Convergence Theorem, $M_n$ converges almost surely
and in $L^1$ to a limit $M_\infty$, such that $M_n = \EE[M_\infty | \mathfrak{F}_n]$
almost surely. Clearly $\abs{M_\infty} \leq K$ almost surely.

Now, the claim is that the indefinite integral of $M_\infty$ is $f$, i.e.
\[
    f(y) - f(0) = \mu\left(M_\infty \id_{x \leq y}\right)
\]
almost everywhere. First let me show this for all dyadic $y = k\cdot 2^{-n}$.
Indeed,
\begin{align*}
    \mu\left(M_\infty \id_{x \leq k\cdot 2^{-n}}\right)
    &= \sum_{i = 0}^{k-1} \mu\left(M_\infty \id_{i \cdot 2^{-n}\leq x < (i+1)\cdot 2^{-n}}\right)\\
    &= \sum_{i = 0}^{k-1} \EE\left[\EE[M_\infty | \mathfrak{F}_n] \id_{i \cdot 2^{-n}\leq x < (i+1)\cdot 2^{-n}}\right]\\
    &= \sum_{i = 0}^{k-1} \EE\left[M_n \id_{i \cdot 2^{-n}\leq x < (i+1)\cdot 2^{-n}}\right]\\
    &= \sum_{i = 0}^{k-1} f((i+1)\cdot 2^{-n}) - f(i\cdot 2^{-n})\\
    &= f(k\cdot 2^{-n}) - f(0)
\end{align*}


For general $y$, we can take a dyadic sequence $y_n$ approaching
$y$ from the left. Then, $\mu(M_\infty \id_{x\leq y}) = \lim \mu(M_\infty \id_{x\leq y_n})$
by dominated convergence (note that $M_\infty$ is bounded), $\mu(M_\infty \id_{x\leq y_n}) = f(y_n)$
for all $y_n$ by our construction, and $\lim f(y_n) = f(y)$ by continuity.

\end{document}